% Conclusion
% viscode_shared_subspace_probe — D028 framing
% Generated: 2026-04-11

\section{Conclusion}
\label{sec:conclusion}

We have presented a training-free probing framework for measuring cross-format representational similarity in visual code language models, combining linear CKA with bootstrap confidence intervals, permutation null testing, and Procrustes distance robustness check.

Applied to six models spanning four architecture families (Qwen, Mistral, BigCode, MoE) across SVG, TikZ, and Asymptote (252 matched triples, 7 equidistant layers per model), we find:
(i)~statistically significant but weak cross-format CKA in all six models ($2.55$--$3.01\times$ the permutation null, $p < 0.001$);
(ii)~a universal TikZ--Asy $\gg$ SVG--TikZ asymmetry consistent with a token-overlap gradient (BPE Jaccard $\rho \approx 0.7$ at pair level);
(iii)~architecture-dependent layer profiles---inverted-U, broadly rising, late-plateau, and bimodal---ruling out a universal ``mid-layer peak'' characterization; and
(iv)~a pre-registered non-visual Python negative control (D060) formally fails: python code yields 60--110\% of the visual CKA across 20/21 evaluation cells, establishing that code-general structure---rather than visual-specific representations---dominates the measured alignment.

The cross-format similarity is of \textbf{uncertain origin}: token-overlap gradient and code-general structure provide plausible non-visual explanations for the majority of the signal.
We release the complete analysis pipeline as a methodological template for future work on format-residualized CKA and visually-paired datasets to isolate genuine visual abstraction.
